\documentclass[10pt]{article}
\usepackage[margin=0.85in]{geometry}
\usepackage{amsmath,amssymb,amsthm}
\usepackage{graphicx}
\usepackage{microtype}
\usepackage[hidelinks]{hyperref}

\newtheorem{theorem}{Theorem}
\newcommand{\Hh}{\mathcal H}
\newcommand{\supp}{\operatorname{supp}}

\title{A nonatomic counterexample to the\\Measure Minimization Conjecture}
\author{Counterexample packet for arXiv:2308.00312}
\date{11 August 2026}

\begin{document}
\maketitle

\noindent\textbf{Status.} Full counterexample to Conjecture 3.2 as stated;
likely valid, pending human review.

\section{Source conjecture}

Krishna's Measure Minimization Conjecture \cite{Krishna} asserts the following.
Let $\{\tau_\alpha\}_{\alpha\in\Omega}$ be a continuous frame for a Hilbert
space $\Hh$. If
\[
 h=\theta_\tau^*f,\qquad f\in L^2(\Omega,\mu),
\]
and
\begin{equation}\label{eq:threshold}
 \mu(\supp f)<\frac12\left(1+
 \frac{1}{\sup_{\alpha\ne\beta}|\langle\tau_\alpha,
 \tau_\beta\rangle|}\right),
\end{equation}
then $f$ is claimed to be the unique minimizer of support measure among all
$g\in L^2(\Omega,\mu)$ satisfying $\theta_\tau^*g=h$.

\begin{center}
\includegraphics[width=0.96\textwidth]{figures/open_problem_crop.png}
\end{center}

\section{Proof intuition}

Cardinality sparsity in a finite dictionary cannot be reduced continuously:
one cannot use half an atom.  Lebesgue support measure behaves differently.
For a constant continuous frame, a coefficient with a prescribed integral can
be compressed into an arbitrarily small interval simply by increasing its
amplitude.  The frame remains Parseval and the conjectured coherence threshold
remains fixed, but the support-measure problem has infimum zero.  The formal
example below implements exactly this obstruction.

\newpage
\section{Counterexample}
\medskip

\begin{theorem}
Conjecture 3.2 of \cite{Krishna} is false for a normalized continuous Parseval
frame on the one-dimensional real Hilbert space.
\end{theorem}

\begin{proof}
Let
\[
 (\Omega,\mu)=([0,1],\lambda),\qquad \Hh=\mathbb R,
 \qquad \tau_\alpha=1\quad(\alpha\in[0,1]).
\]
This is a normalized continuous Parseval frame: for every $h\in\mathbb R$,
\[
 \int_0^1|\langle h,\tau_\alpha\rangle|^2\,d\alpha
 =\int_0^1|h|^2\,d\alpha=|h|^2.
\]
Its synthesis operator is
\[
 \theta_\tau^*f=\int_0^1 f(\alpha)\tau_\alpha\,d\alpha
 =\int_0^1 f(\alpha)\,d\alpha,
\]
and its off-diagonal coherence is
\[
 \sup_{\alpha\ne\beta}|\langle\tau_\alpha,\tau_\beta\rangle|=1.
\]

Fix any $a\in(0,1)$ and set
\[
 f_a=\frac1a\mathbf 1_{(0,a)}\in L^2([0,1]).
\]
Then
\[
 \theta_\tau^*f_a=\frac1a\lambda(0,a)=1,
 \qquad \mu(\supp f_a)=a<1
 =\frac12\left(1+\frac11\right).
\]
Thus $f_a$ satisfies every hypothesis of the conjecture for $h=1$.

Now choose any $b\in(0,a)$ and define
$f_b=b^{-1}\mathbf1_{(0,b)}$.  Again $f_b\in L^2([0,1])$ and
\[
 \theta_\tau^*f_b=1=\theta_\tau^*f_a,
 \qquad \mu(\supp f_b)=b<a=\mu(\supp f_a).
\]
Therefore $f_a$ is not even a minimizer of the problem $(P_M)$, and hence
cannot be its unique minimizer.  Since $b$ can tend to zero, the infimum of
the support measure in the synthesis fiber of $h=1$ is zero; it is not
attained because the zero-support equivalence class synthesizes zero.
\end{proof}

\section{Verification, scope, and novelty}

The accompanying exact-arithmetic script checks the frame identity, coherence,
threshold, synthesis values, and strict support decrease for several rational
pairs $0<b<a<1$.  These checks are illustrative; the displayed calculation is
the proof.

The counterexample resolves only Conjecture 3.2.  It does not resolve the
separate questions about equality, endpoint versions, and Tao analogues, or
the NP-hardness classification in Problem 3.3.  The obstruction suggests that a repaired
statement must impose atomicity or amplitude control, or replace support
measure by a coercive objective such as total variation.

A bounded search on 11 August 2026 used the exact conjecture name and formula,
the source id and author, and continuous-frame support-minimization phrases.
Later related preprints by the same author on continuous Deutsch, spark, and
Welch inequalities were checked at title/abstract level. No prior
counterexample or retraction was located. Novelty confidence is moderate
because the example is elementary.

\paragraph{Human-review recommendation.}
Verify the standard definition of a continuous Parseval frame and the
one-line synthesis calculation.  If confirmed, classify this as a full
counterexample to Conjecture 3.2 as literally stated.

\begin{thebibliography}{9}
\bibitem{Krishna}
K. M. Krishna,
\emph{Functional Continuous Uncertainty Principle},
arXiv:2308.00312 (2023), Conjecture 3.2.
\end{thebibliography}

\end{document}
